\section{Checking an apparent saving}
\label{sec:workplace}

Before randomisation, each team fixes its set of tasks and scores their requirements. Suppose a team takes the two export tasks and both pass their acceptance checks by the cutoff. At follow-up the March export contains both orders. Including the April order costs the filter one unit. The repaired export still prints order number 007 with both leading zeros, so the repair retains its unit. The team retains two of its three initial units ($Q/S=2/3$). Its labour per unit includes all hours on those requirements through age $\tau+\Delta$, including work after the illustrated shift.

An independent fair draw in each matched pair selects the team offered the specified AI tools. For hours per initial unit and for the retained fraction, subtract the comparison team's value from the value for the team offered AI. If each team's outcomes depend only on its own offer, averaging these differences over pairs and possible draws gives the mean effect of the offer on the enrolled teams. A shared reviewer or code passed between teams can let one team's offer affect another team's outcome \cite{hudgens2008}.

The complete-data rule requires both outcomes from every team. To declare a benefit, the trial must establish a saving above one person-hour per initial unit with a retention loss below five percentage points. The teams must justify these thresholds before the trial \cite{piaggio2012}. In a separate example with six hypothetical pairs, counting every hour of work and scoring each check by whether it passed puts the mean saving and retention loss exactly at those thresholds. A benefit requires a larger saving and a smaller loss. Omitting one review hour per initial unit from the AI teams' accounts doubles their apparent saving. Each requirement has 100 equally weighted checks. Counting six failures as passes on each AI-team requirement turns a five-point loss into an apparent one-point gain. Together the errors produce a false benefit claim.

To check for omitted hours, select shifts independently of whether contributors replied. Compare their accounts with separate evidence of who worked on the requirements and for how long, including external review within the labour window. Allow for work absent from both sources by bounding how long each contributor could have spent on the requirements during the shift. On these shifts, compare the time contributors spend answering requests about sampled intervals with the time they spend giving a full daily account. Use the independent evidence to check omissions in each method. At follow-up, have independent reviewers retest the acceptance criteria and resolve discrepancies where the evidence permits.

Set the rules for bounding unobserved outcomes before the trial begins. Denote the number of enrolled pairs by $P\ge 1$ and index them by $h$. Calculate a separate difference $D_{ho}$ in each pair for each outcome $o$. Use $o=\mathrm{lab}$ for hours per initial unit and $o=\mathrm{ret}$ for the retained fraction. Before the draw, specify a finite range $[a_{ho},b_{ho}]$ that contains the difference whichever team receives the offer. After the draw, use the available measurements to bound the difference under the allocation actually drawn by $[\ell_{ho},u_{ho}]$. Use the range fixed before the draw if both teams lack that outcome.

Across pairs, the mean difference under the draw is $\overline D_o=P^{-1}\sum_{h=1}^{P}D_{ho}$. Write $\mu_o=\E[\overline D_o]$ for its expectation over possible draws. Fix the teams' potential outcomes under either offer. If outcomes in a pair depend only on that pair's draw, the independent draws make the pair differences independent. Hoeffding's inequality bounds the chance that $\overline D_o$ falls more than $\varepsilon_o$ below $\mu_o$ by $\alpha_o$. The same bound holds above $\mu_o$ \cite{hoeffding1963}. For $0<\alpha_o<1$, the margin is
\[
\varepsilon_o=\sqrt{\frac{\log(1/\alpha_o)}{2P^2}
  \sum_{h=1}^{P}(b_{ho}-a_{ho})^2}.
\]
Let $\overline{\ell}_o$ and $\overline{u}_o$ be the averages of the lower and upper limits across pairs after the draw. If each range contains its pair's difference, then $\overline{\ell}_o\le\overline D_o\le\overline{u}_o$. A benefit requires the upper confidence limit for labour, $\overline{u}_{\rm lab}+\varepsilon_{\rm lab}$, to fall below $-1$ person-hour per initial unit. The lower confidence limit for retention, $\overline{\ell}_{\rm ret}-\varepsilon_{\rm ret}$, must exceed $-0.05$.

For outcome $o$, suppose the chance that any pair's range excludes its actual difference is at most $\gamma_o\in[0,1]$. A confidence limit can exclude $\mu_o$ only if a pair's range excludes its difference or the draw takes the mean beyond the corresponding one-sided margin. By the union bound the chance is at most $\alpha_o+\gamma_o$. If a true effect fails its threshold, a false benefit claim requires the confidence limit for that outcome to fail. The larger of the two error bounds therefore covers the chance of a false claim whichever threshold fails.

For 48 pairs with labour differences between $-10$ and $10$ person-hours per initial unit, $\alpha_{\rm lab}=0.025$ gives a margin of about 3.92 in those units. Even with exact measurements, the observed mean saving would have to exceed about 4.92 person-hours per initial unit to establish a saving above one. On the preliminary shifts, check which contributors could have worked on each requirement and how many hours each had available. Use those hours and each team's initial score to justify the range of possible differences in labour per initial unit. Before the main trial, simulate the rule across effects that meet or fail its thresholds to estimate how often it detects a benefit or declares one in error.
